\documentclass[a4paper, 12pt]{extreport}


\usepackage{fontenc}
\usepackage{tempora}
\usepackage[utf8]{inputenc}
\usepackage[english,russian]{babel}
\usepackage{amssymb,amsfonts,amsmath,mathtext,cite,enumerate,float}
\usepackage{pgfplots}
\usepackage{graphicx}
\usepackage{tocloft}
\usepackage{listings}
\usepackage[tableposition=top]{caption}
\usepackage{titlesec}
\usepackage{setspace}
\usepackage{geometry}
\usepackage{indentfirst}
\usepackage{pdfpages}
\usepackage{enumerate,letltxmacro}
\usepackage{threeparttable}
\usepackage[hidelinks]{hyperref}
\usepackage{flafter}
\usepackage{enumitem}
\usepackage{multirow}
\usepackage[figure,table]{totalcount}
\usepackage{lastpage}
\usepackage{longtable}
\usepackage{array}
\usepackage[most]{tcolorbox}
\usepackage{xcolor}
\usepackage{xparse}
\usepackage{textcomp}


\setlist{nosep}

\setcounter{secnumdepth}{3}


\newcommand{\ssr}[1]{
	\begin{center}
		\LARGE\bfseries{#1}
	\end{center} 
	\addcontentsline{toc}{chapter}{#1}  
}

\newcommand{\ssrn}[1]{%
  \begin{center}
    \LARGE\bfseries #1
  \end{center}%
   \vspace{1em}
}





\DeclareUnicodeCharacter{2713}{}  % ✓
\DeclareUnicodeCharacter{2717}{X} % ✗
\DeclareUnicodeCharacter{274C}{X} % ❌
\DeclareUnicodeCharacter{1F4E6}{} % 📦
\DeclareUnicodeCharacter{2705}{}  % ✅
\DeclareUnicodeCharacter{1F50D}{} % 🔍
\DeclareUnicodeCharacter{1F4CC}{} % 📌
\DeclareUnicodeCharacter{200A}{~}
\DeclareUnicodeCharacter{202F}{~}





\newcommand{\russianAppendixLetter}[1]{%
  \ifcase#1%
  \or А\or Б\or В\or Г\or Д\or Е\or Ж\or И\or К\or Л\or М\or Н\or П\or Р\or С\or Т\or У\or Ф\or Х\or Ц\or Ш\or Щ\or Э\or Ю\or Я%
  \else\@ctrerr % Ошибка, если счётчик вышел за пределы доступных букв
  \fi
}

% Команда для начала раздела приложений
\newcommand{\BeginAppendices}{%
  \clearpage % Начать приложения с новой страницы
  \appendix % Стандартная команда LaTeX для переключения в режим приложений
            % Она сбрасывает счётчик глав и меняет \thechapter на \Alph{chapter}

  % Переопределяем отображение номера главы на нашу русскую букву
  % \appendix уже сбросил chapter в 0. \refstepcounter{chapter} в \prilojenie сделает его 1 для первого приложения.
  \renewcommand{\thechapter}{\russianAppendixLetter{\value{chapter}}}

  % Меняем название "Глава" на "Приложение" для заголовков и TOC
  % Для классов report и book это \chaptername. Но \appendix может уже менять \chaptername на \appendixname.
  % \renewcommand{\chaptername}{Приложение} % Это изменит "Глава X" на "Приложение X"
  % Однако, мы хотим "Приложение А", "Приложение Б" в заголовке, а не "Приложение А Название"
  % Поэтому мы будем форматировать заголовок полностью в команде \prilojenie.
  % Но для TOC это может быть полезно, если стиль TOC использует \chaptername.
  % Вместо этого, \addcontentsline будет формировать нужную строку.

  % Переопределяем нумерацию для секций, рисунков, таблиц, листингов, чтобы они включали букву приложения
  \renewcommand{\thesection}{\thechapter.\arabic{section}}
  \renewcommand{\thesubsection}{\thesection.\arabic{subsection}}
  \renewcommand{\thefigure}{\thechapter.\arabic{figure}}
  \renewcommand{\thetable}{\thechapter.\arabic{table}}
  \renewcommand{\thelstlisting}{\thechapter.\arabic{lstlisting}} % КЛЮЧЕВОЕ для листингов

  % Убедимся, что имя для листинга остаётся "Листинг" (буква А будет добавлена через \thelstlisting)
  \renewcommand{\lstlistingname}{Листинг}

  % Можно добавить общий заголовок "Приложения" в TOC (необязательно)
  % \addcontentsline{toc}{part}{Приложения} % Как "Часть"
  % или использовать \part*{Приложения} для заголовка на отдельной странице
}

% Ваша команда \prilojenie, адаптированная
% Принимает один аргумент: название приложения (например, "Листинг GUI клиента")
\newcommand{\prilojenie}[1]{%
  % \chapter{#1} % Стандартная команда \chapter теперь будет работать как надо,
                % но она выведет "Приложение А: Название приложения" или похоже,
                % в зависимости от стиля. Нам нужно customnoe форматирование.

  \clearpage % Каждое приложение с новой страницы (если \chapter этого не делает в вашем стиле)
  \refstepcounter{chapter} % Увеличиваем счётчик глав (теперь это наш счётчик приложений А, Б, В...)
                           % \refstepcounter также делает метку для \label/\ref
                           % Важно: \appendix сбрасывает chapter в 0, так что первое приложение будет "1" -> "А"

  % Форматирование заголовка приложения
  \phantomsection % Для корректной работы ссылок hyperref
  % \addtocontents{toc}{\protect\addvspace{10pt}} % Добавить немного места в TOC перед записями приложений (если нужно)
  
  \begin{center}
    \LARGE\bfseries Приложение \thechapter % \thechapter теперь наша русская буква
  \end{center}
  \noindent\textbf{#1}
  
  % Добавляем в содержание (Table of Contents)
  % Формат: "Приложение А. Название приложения"
  \addcontentsline{toc}{chapter}{Приложение \thechapter. #1}
  
  % Установка заголовков для страниц (headers/footers)
  \markboth{\MakeUppercase{Приложение \thechapter. #1}}{\MakeUppercase{Приложение \thechapter. #1}}

  % После \refstepcounter{chapter}, счётчики, зависимые от главы (секции, рисунки, таблицы, листинги),
  % должны автоматически сброситься в 0, если они так настроены (обычно так и есть).
  % Например, счётчик lstlisting сбрасывается при каждом \chapter.
}

\newcommand{\abzend}{\newline\indent}

\makeatletter
\renewcommand\LARGE{\@setfontsize\LARGE{22pt}{20}}
\renewcommand\Large{\@setfontsize\Large{20pt}{20}}
\renewcommand\large{\@setfontsize\large{16pt}{20}}
\makeatother
\RequirePackage{titlesec}
\titleformat{\chapter}[block]{\large\bfseries}{\thechapter}{0.5em}{\large\bfseries}
\titleformat{name=\chapter,numberless}[block]{\hspace{\parindent}}{}{0pt}{\large\bfseries\centering}
\titleformat{\section}[block]{\large\bfseries}{\thesection}{0.5em}{\large\bfseries}
\titleformat{\subsection}[block]{\large\bfseries}{\thesubsection}{0.5em}{\large\bfseries}
\titleformat{\subsubsection}[block]{\large\bfseries}{\thesubsubsection}{0.5em}{\large\bfseries}
\titlespacing{\chapter}{\parindent}{-22pt}{10pt}
\titlespacing{\section}{\parindent}{10pt}{10pt}
\titlespacing{\subsection}{\parindent}{10pt}{10pt}
\titlespacing{\subsubsection}{\parindent}{10pt}{10pt}
\titlespacing{\paragraph}{\parindent}{10pt}{10pt}

\makeatletter
\renewcommand{\@biblabel}[1]{#1.}
\makeatother
%
%\titleformat{\chapter}[hang]{\LARGE\bfseries}{\hspace{1.25cm}\thechapter}{1ex}{\LARGE\bfseries}
%\titleformat{\section}[hang]{\Large\bfseries}{\hspace{1.25cm}\thesection}{1ex}{\Large\bfseries}
%\titleformat{name=\section,numberless}[hang]{\Large\bfseries}{\hspace{1.25cm}}{0pt}{\Large\bfseries}
%\titleformat{\subsection}[hang]{\large\bfseries}{\hspace{1.25cm}\thesubsection}{1ex}{\large\bfseries}
%\titlespacing{\chapter}{0pt}{-\baselineskip}{\baselineskip}
%\titlespacing*{\section}{0pt}{\baselineskip}{\baselineskip}
%\titlespacing*{\subsection}{0pt}{\baselineskip}{\baselineskip}

\geometry{left=31mm}
\geometry{right=14mm}
\geometry{top=20mm}
\geometry{bottom=20mm}

\onehalfspacing

\renewcommand{\theenumi}{\arabic{enumi}}
\renewcommand{\labelenumi}{\arabic{enumi}\text{)}}
\renewcommand{\theenumii}{.\arabic{enumii}}
\renewcommand{\labelenumii}{\arabic{enumii}\text{)}}
\renewcommand{\theenumiii}{.\arabic{enumiii}}
\renewcommand{\labelenumiii}{\arabic{enumi}.\arabic{enumii}.\arabic{enumiii}.}

\renewcommand{\cftchapleader}{\cftdotfill{\cftdotsep}}

%\addto\captionsrussian{\renewcommand{\figurename}{Рисунок}}
%\DeclareCaptionLabelSeparator{dash}{~---~}
%\captionsetup{labelsep=dash}

%\captionsetup[figure]{justification=centering,labelsep=dash}
%\captionsetup[table]{labelsep=dash,justification=raggedright,singlelinecheck=off}


% НАСТРОЙКА РИСУНКОВ И ТАБЛИЦ

\usepackage{subcaption}
\DeclareCaptionLabelFormat{gostfigure}{Рисунок #2}
\DeclareCaptionLabelFormat{gosttable}{Таблица #2}
\DeclareCaptionLabelSeparator{gost}{~---~}
\captionsetup{labelsep=gost}
\captionsetup[figure]{labelformat=gostfigure}
\captionsetup[table]{labelformat=gosttable,justification=raggedright,singlelinecheck=off}
\renewcommand{\thesubfigure}{\asbuk{subfigure}}




\graphicspath{{images/}{../code/res}}%путь к рисункам

\newcommand{\floor}[1]{\lfloor #1 \rfloor}


\lstdefinestyle{basestyle}{ %
	basicstyle=\ttfamily\small, % размер и начертание шрифта для подсветки кода
	upquote=true,
	numbers=left,               % где поставить нумерацию строк (слева\справа)
	numberstyle=\tiny,           % размер шрифта для номеров строк
	stepnumber=1,                   % размер шага между двумя номерами строк
	numbersep=7pt,                % как далеко отстоят номера строк от подсвечиваемого кода
	xleftmargin=15pt,        % Отступ слева, чтобы номера строк были внутри рамки
    framexleftmargin=15pt,    % Расширение рамки слева, чтобы номера строк были внутри
	showspaces=false,            % показывать или нет пробелы специальными отступами
	showstringspaces=false,      % показывать или нет пробелы в строках
	showtabs=false,             % показывать или нет табуляцию в строках
%	frame=single,              % рисовать рамку вокруг кода
	tabsize=4,                 % размер табуляции по умолчанию равен 2 пробелам
	captionpos=t,              % позиция заголовка вверху [t] или внизу [b] 
	breaklines=true,           % автоматически переносить строки (да\нет)
	breakatwhitespace=false, % переносить строки только если есть пробел
	escapeinside={\#*}{*)},   % если нужно добавить комментарии в коде
	abovecaptionskip=-5pt,
	inputencoding=utf8, 
	extendedchars=\true,      % вроде было нужно для отображения ru символов
	keepspaces=true,           % вроде было нужно для отображения ru символов
	literate={«}{{\guillemotleft}}1            % Кавычка-ёлочка левая
             {»}{{\guillemotright}}1           % Кавычка-ёлочка правая
             {…}{{\dots}}1                     % Многоточие (как один символ Unicode)
             {—}{{\textemdash}}1               % Длинное тире (Em-dash)
             {–}{{\textendash}}1               % Короткое тире (En-dash)
             {→}{{$\rightarrow$}}1             % Стрелка вправо
             {←}{{$\leftarrow$}}1              % Стрелка влево
             {а}{{\char224}}1 {б}{{\char225}}1 {в}{{\char226}}1 {г}{{\char227}}1 {д}{{\char228}}1 {е}{{\char229}}1 {ё}{{\"e}}1 {ж}{{\char230}}1 {з}{{\char231}}1 {и}{{\char232}}1 {й}{{\char233}}1 {к}{{\char234}}1 {л}{{\char235}}1 {м}{{\char236}}1 {н}{{\char237}}1 {о}{{\char238}}1 {п}{{\char239}}1 {р}{{\char240}}1 {с}{{\char241}}1 {т}{{\char242}}1 {у}{{\char243}}1 {ф}{{\char244}}1 {х}{{\char245}}1 {ц}{{\char246}}1 {ч}{{\char247}}1 {ш}{{\char248}}1 {щ}{{\char249}}1 {ъ}{{\char250}}1 {ы}{{\char251}}1 {ь}{{\char252}}1 {э}{{\char253}}1 {ю}{{\char254}}1 {я}{{\char255}}1 {А}{{\char192}}1 {Б}{{\char193}}1 {В}{{\char194}}1 {Г}{{\char195}}1 {Д}{{\char196}}1 {Е}{{\char197}}1 {Ё}{{\"E}}1 {Ж}{{\char198}}1 {З}{{\char199}}1 {И}{{\char200}}1 {Й}{{\char201}}1 {К}{{\char202}}1 {Л}{{\char203}}1 {М}{{\char204}}1 {Н}{{\char205}}1 {О}{{\char206}}1 {П}{{\char207}}1 {Р}{{\char208}}1 {С}{{\char209}}1 {Т}{{\char210}}1 {У}{{\char211}}1 {Ф}{{\char212}}1 {Х}{{\char213}}1 {Ц}{{\char214}}1 {Ч}{{\char215}}1 {Ш}{{\char216}}1 {Щ}{{\char217}}1 {Ъ}{{\char218}}1 {Ы}{{\char219}}1 {Ь}{{\char220}}1 {Э}{{\char221}}1 {Ю}{{\char222}}1 {Я}{{\char223}}1
}

\lstdefinestyle{pythonstyle}{
    style=basestyle,
    language=Python
}

\lstdefinestyle{logstyle}{
    style=basestyle,
    language={},                % Убираем язык, чтобы не подсвечивало синтаксис
    columns=fullflexible,       % Лучшее выравнивание для текстовых логов
    commentstyle=\color{black}, % На всякий случай сбрасываем цвета
    stringstyle=\color{black},
    keywordstyle=\color{black}
}

\lstset{style=basestyle}

% --- 2. Логика на LaTeX3 ---
\ExplSyntaxOn

% Переменные
\tl_new:N \l_mycode_caption_tl
\tl_new:N \l_mycode_label_tl
\tl_new:N \l_mycode_params_tl

\keys_define:nn { mycode } {
    caption .tl_set:N = \l_mycode_caption_tl,
    label .tl_set:N = \l_mycode_label_tl,
    
    unknown .code:n = {
        \tl_if_empty:nTF { #1 }
        { \tl_put_right:Nx \l_mycode_params_tl { , \l_keys_key_tl } }
        { 
            \tl_put_right:Nx \l_mycode_params_tl { , \l_keys_key_tl = }
            \tl_put_right:Nn \l_mycode_params_tl { #1 }
        }
    }
}


% --- 3. Команда \inputcode ---
\NewDocumentCommand{\inputcode}{ +O{} m }{
    % Очистка
    \tl_clear:N \l_mycode_caption_tl
    \tl_clear:N \l_mycode_label_tl
    \tl_clear:N \l_mycode_params_tl
    \tl_set:Nn \l_mycode_params_tl { style=logstyle, aboveskip=0pt, belowskip=0pt }
    
    % Парсинг
    \keys_set:nn { mycode } { #1 }
    
    % Увеличиваем счётчик и ставим label
    \refstepcounter{lstlisting}
    \tl_if_empty:NF \l_mycode_label_tl { \label{\l_mycode_label_tl} }
    
    % --- ГЕНЕРАЦИЯ ЗАГОЛОВКОВ (В ОДНОМ МЕСТЕ) ---
    % 1. Формируем основной заголовок: "Листинг X.Y --- Название"
    \tl_set:Nn \l_mycode_title_main_tl {
        Листинг~\thelstlisting
        \tl_if_empty:NTF \l_mycode_caption_tl {} {~---~\l_mycode_caption_tl }
    }
    
    % 2. Формируем заголовок продолжения"
    \tl_set:Nn \l_mycode_title_cont_tl {
        Листинг~\thelstlisting
        \tl_if_empty:NTF \l_mycode_caption_tl {} {~(продолжение)~---~\l_mycode_caption_tl }
    }
    
    % Рисуем блок
    \tcbinputlisting{
        listing~file={#2},
        listing~only,
        listing~options/.expand~once = \l_mycode_params_tl,
        
        enhanced, breakable, colback=white, colframe=white, boxrule=0pt,
        sharp~corners, left=0pt, right=0pt, top=0pt, bottom=0pt,
        before~skip=20pt,
        
        
        % 1. Неразорванный (Основной заголовок)
        overlay~unbroken={
            \draw[black, line~width=0.5pt] (frame.north~west) rectangle (frame.south~east);
            \node[anchor=south~west, yshift=2pt, inner~sep=0pt] at (frame.north~west) {
                \l_mycode_title_main_tl
            };
        },
        
        % 2. Начало (Основной заголовок)
        overlay~first={
            \draw[black, line~width=0.5pt] (frame.north~west) rectangle (frame.south~east);
            \node[anchor=south~west, yshift=2pt, inner~sep=0pt] at (frame.north~west) {
                \l_mycode_title_main_tl
            };
        },
        
        % 3. Середина (Продолжение)
        overlay~middle={
            \draw[black, line~width=0.5pt] (frame.north~west) rectangle (frame.south~east);
            \node[anchor=south~west, yshift=2pt, inner~sep=0pt] at (frame.north~west) {
                \l_mycode_title_cont_tl
            };
        },
        
        % 4. Конец (Продолжение)
        overlay~last={
            \draw[black, line~width=0.5pt] (frame.north~west) rectangle (frame.south~east);
            \node[anchor=south~west, yshift=2pt, inner~sep=0pt] at (frame.north~west) {
                \l_mycode_title_cont_tl
            };
        }
    }
}
\ExplSyntaxOff

\captionsetup[lstlisting]{singlelinecheck=false, justification=raggedright}

\pgfplotsset{width=0.85\linewidth, height=0.5\columnwidth}

\linespread{1.3}

\parindent=1.25cm

%\LetLtxMacro\itemold\item
%\renewcommand{\item}{\itemindent0.75cm\itemold}
% label=\textbullet
\def\labelitemi{---}
\setlist[itemize]{leftmargin=0.65cm, itemindent=0.65cm}
\setlist[enumerate]{leftmargin=0.65cm, itemindent=0.55cm}

\newcommand{\specialcell}[2][c]{%
	\begin{tabular}[#1]{@{}c@{}}#2\end{tabular}}

\renewcommand{\url}[1]{\href{#1}{\detokenize{#1}}}

\newcommand{\n}{\newline}

\frenchspacing
